\documentclass[11pt,a4paper]{article}
\usepackage[T1]{fontenc}
\usepackage[utf8]{inputenc}
\usepackage{lmodern}
\usepackage{amsmath,amssymb,amsthm,mathtools}
\usepackage{microtype}
\usepackage[margin=27mm]{geometry}
\usepackage{enumitem}
\usepackage{xcolor}
\usepackage[colorlinks=true,linkcolor=blue!50!black,citecolor=blue!50!black,urlcolor=blue!50!black]{hyperref}
\hypersetup{pdftitle={Prime-Logarithm Integrality and Perfect-Power Obstructions},pdfauthor={Keio},pdfsubject={Periodic harmonic sums, ordinary distributions, and local integrality}}
\setlength{\emergencystretch}{2em}
\setlist[enumerate]{itemsep=3pt,topsep=5pt}
\newtheorem{theorem}{Theorem}[section]
\newtheorem{proposition}[theorem]{Proposition}
\newtheorem{lemma}[theorem]{Lemma}
\newtheorem{corollary}[theorem]{Corollary}
\theoremstyle{definition}
\newtheorem{definition}[theorem]{Definition}
\newtheorem{example}[theorem]{Example}
\theoremstyle{remark}
\newtheorem{remark}[theorem]{Remark}
\newcommand{\Q}{\mathbb Q}
\newcommand{\Z}{\mathbb Z}
\newcommand{\F}{\mathbb F}
\newcommand{\C}{\mathbb C}
\newcommand{\Zloc}{\Z_{(\ell)}}
\DeclareMathOperator{\spanop}{span}
\DeclareMathOperator{\rank}{rank}
\DeclareMathOperator{\ord}{ord}
\DeclareMathOperator{\im}{im}
\title{Prime-Logarithm Integrality in Digamma Lattices\\
and Perfect-Power Obstructions for Periodic Harmonic Sums}
\author{Keio}
\date{15 November 2026\\\small Research manuscript}

\begin{document}
\maketitle

\begin{abstract}
Let $\ell$ be a prime and let $q>1$ have all its prime divisors congruent
to $1$ modulo $\ell$. We prove that the $\Z_{(\ell)}$-span of the numbers
$\psi(a/q)+\gamma$, $1\le a<q$, intersects the rational span of the prime
logarithms $\log p$, $p\mid q$, in exactly their $\Z_{(\ell)}$-span.
The proof compares the rational dimension of the digamma span, determined
by Murty--Pathak, with the rank in characteristic $\ell$ of the ordinary
distribution relations, determined by Kubert's integral freeness theorem.
Equality of these ranks gives a saturated presentation and independence
of the prime logarithms modulo $\ell$. As an application, there exists a
mean-zero, $q$-periodic integer function with $f(0)=0$, with
$f(a)\equiv1\pmod\ell$ for $1\le a<q$, and with
$\sum_{n\ge1}f(n)/n=0$ if and only if $q$ is a perfect $\ell$th power.
For $\ell=2$, this classifies all odd periods admitting odd-valued
integer harmonic relations. In particular, every convergent Erd\H{o}s
sign sum with a nonsquare odd period is transcendental. We also prove
the corresponding necessary condition for coefficients congruent to
$1$ modulo $\ell^k$ and give explicit constructions and examples.
\end{abstract}

\noindent\textbf{Status.} The arguments are presented as research proofs.
External peer review and discovery priority have not been established.

\section{Statements and conventions}

Throughout, $q>1$ is an integer, $\omega(q)$ is the number of its distinct
prime divisors, and $\varphi(q)$ is Euler's totient. For a prime $\ell$,
write
\[
 \Zloc=\left\{\frac ab\in\Q:\ell\nmid b\right\}.
\]
Congruences between elements of $\Zloc$ are taken in this local ring.
The notation $v_p(q)$ denotes the exponent of $p$ in $q$.

Let $\psi=\Gamma'/\Gamma$ be the digamma function and let $\gamma$ be
Euler's constant. Put
\begin{equation}\label{eq:digamma-values}
 D_a(q)=\psi(a/q)+\gamma,\qquad 1\le a<q,
 \qquad
 \Lambda_{q,\ell}
 =\sum_{a=1}^{q-1}\Zloc D_a(q).
\end{equation}
All logarithms of positive real numbers in this paper are real
logarithms. Our structural result concerns the integral coefficients of
relations among the values in \eqref{eq:digamma-values}.

\begin{theorem}[Prime-logarithm integrality]\label{thm:purity-main}
Let $\ell$ be prime and assume that
\begin{equation}\label{eq:residue-hypothesis}
 p\mid q\quad\Longrightarrow\quad p\equiv1\pmod\ell.
\end{equation}
Then
\begin{equation}\label{eq:purity-main}
 \Lambda_{q,\ell}\cap
 \spanop_{\Q}\{\log p:p\mid q\}
 =\bigoplus_{p\mid q}\Zloc\log p.
\end{equation}
Moreover, the classes of $\log p$, $p\mid q$, are linearly independent
over $\F_\ell$ in $\Lambda_{q,\ell}/\ell\Lambda_{q,\ell}$.
\end{theorem}

For a rational-valued periodic function $f$, we use the notation
\[
 L(1,f)=\sum_{n=1}^{\infty}\frac{f(n)}n
\]
only when this ordinary series converges. Convergence is equivalent to
the zero-mean condition over a period. The next theorem is an exact
existence classification, including an integer-valued converse.

\begin{theorem}[Perfect-power classification]\label{thm:classification}
Assume \eqref{eq:residue-hypothesis}. The following are equivalent.
\begin{enumerate}[label=\textup{(\roman*)}]
\item There is a $q$-periodic rational-valued function $f$ such that
\begin{equation}\label{eq:function-conditions}
 f(0)=0,\qquad
 f(a)\in1+\ell\Zloc\ (1\le a<q),\qquad
 \sum_{a=1}^{q}f(a)=0,\qquad L(1,f)=0.
\end{equation}
\item There is a function as in \textup{(i)} whose values are integers
and satisfy $f(a)\equiv1\pmod\ell$ for $1\le a<q$.
\item The integer $q$ is a perfect $\ell$th power.
\end{enumerate}
No assertion that $q$ is the least period is required.
\end{theorem}

Every odd prime is $1$ modulo $2$, so the case $\ell=2$ has no further
restriction on the odd integer $q$.

\begin{corollary}[Square-period classification]\label{cor:squares}
Let $q>1$ be odd. A mean-zero, $q$-periodic function $f$ with $f(0)=0$
and odd integer values at every other residue class can satisfy
$L(1,f)=0$ if and only if $q$ is a square. The same equivalence holds
when ``odd integer'' is replaced by ``odd element of $\Z_{(2)}$''.
\end{corollary}

\begin{corollary}[Nonsquare Erd\H{o}s sign sums]\label{cor:erdos}
Let $q>1$ be odd and nonsquare. Suppose that $f$ is $q$-periodic,
$f(0)=0$, and $f(a)\in\{-1,1\}$ for $1\le a<q$. If
$\sum_{a=1}^{q}f(a)=0$, then $L(1,f)$ is transcendental.
\end{corollary}

The sign condition in Corollary~\ref{cor:erdos} is stronger than the
congruence condition in Corollary~\ref{cor:squares}. In particular,
the square-period existence statement does not construct a vanishing
sign sum at every square. The explicit witnesses below generally have
values other than $\pm1$.

Corollary~\ref{cor:squares} resolves the square-period conjecture
formulated in VMRC manuscript~[055, \S5.2]~\cite{VMRC055}: an odd
period admits an odd Okada vector exactly when it is a square.
That manuscript reports exact nonexistence checks for odd nonsquares
$q\le13105$. The present proof is uniform in $q$, with its structural
inputs taken directly from the established results below.

The proof uses two established results, recalled precisely in the next
section. The new deduction is a comparison of the same relation matrix
in characteristic zero and characteristic $\ell$. It does not require
an assumption about the arithmetic nature of $\gamma$.

\section{The two structural inputs}

\subsection{The rational digamma span}

Let $\mathcal F^{(0)}(q)$ be the rational vector space of mean-zero,
$q$-periodic functions, with no restriction on $f(0)$. Let
\[
 \mathcal O(q)=\{f\in\mathcal F^{(0)}(q):L(1,f)=0\}.
\]
Murty and Pathak determine an explicit basis of this Okada space and
prove
\begin{equation}\label{eq:okada-dimension}
 \dim_\Q\mathcal O(q)=q-\varphi(q)-\omega(q)
\end{equation}
\cite[\S2 and Theorem~4.4]{MurtyPathak}. The dimension is zero when
$q$ is prime. The allowance of an arbitrary value at the zero residue
class is part of this definition.

For every mean-zero periodic function, the standard digamma expression
is
\begin{equation}\label{eq:digamma-L}
 L(1,f)=-\frac1q\sum_{a=1}^{q}f(a)\psi(a/q).
\end{equation}
Given arbitrary rational numbers $c_a$, $1\le a<q$, extend them by
$f(a)=c_a$ and $f(q)=-\sum_{a<q}c_a$. Since $\psi(1)=-\gamma$,
formula~\eqref{eq:digamma-L} becomes
\[
 L(1,f)=-\frac1q\sum_{a=1}^{q-1}c_aD_a(q).
\]
Thus rank--nullity and \eqref{eq:okada-dimension} give the exact input
we need:
\begin{equation}\label{eq:digamma-dimension}
 \dim_\Q\spanop_\Q\{D_a(q):1\le a<q\}
 =\varphi(q)+\omega(q)-1.
\end{equation}
This is the dimension of the space denoted $D^{(0)}(q)$ by
Murty--Pathak. It must not be confused with a dimension formula for
the smaller space in which $f(0)=0$ is imposed in advance.

We will also use the digamma multiplication formula
\begin{equation}\label{eq:multiplication}
 \sum_{j=0}^{m-1}\psi\left(z+\frac jm\right)
 =m\psi(mz)-m\log m.
\end{equation}
Taking $z=1/m$ and using $\psi(1)=-\gamma$ yields
\begin{equation}\label{eq:total-digamma}
 \sum_{a=1}^{q-1}D_a(q)=-q\log q.
\end{equation}
For each $p\mid q$, the same identity at level $p$ gives
\begin{equation}\label{eq:prime-log-in-span}
 \sum_{j=1}^{p-1}D_{jq/p}(q)=-p\log p.
\end{equation}
Under \eqref{eq:residue-hypothesis}, $p$ is a unit in $\Zloc$, so
$\log p\in\Lambda_{q,\ell}$.

\subsection{Finite-level ordinary distributions}

Let $A_q$ be the free abelian group on symbols $[x]$ indexed by
$x\in q^{-1}\Z/\Z$. Its ordinary distribution subgroup is generated by
\begin{equation}\label{eq:ordinary-relations}
 \sum_{py=x}[y]-[x],
 \qquad p\mid q,\qquad x\in(q/p)^{-1}\Z/\Z.
\end{equation}
All preimages in \eqref{eq:ordinary-relations} belong to $q^{-1}\Z/\Z$.
Kubert's finite-level theorem states that
\begin{equation}\label{eq:kubert}
 U(q)=A_q/\langle\text{relations \eqref{eq:ordinary-relations}}\rangle
 \quad\text{is free abelian of rank }\varphi(q).
\end{equation}
See \cite[Theorems~1.8 and~4.14]{Kubert}. The finite-level presentation,
including its agreement with the relations inherited from the universal
group, is important here. Anderson gives the same statement as a
special case of a weighted freeness theorem
\cite[Theorem~4.3.1(I) and~\S5.3]{Anderson}.

Kubert's presentation may equivalently use multiplication by every
divisor of $q$. Prime multiplication generates those relations:
if $r_{m,x}=\sum_{my=x}[y]-[x]$, then
\[
 r_{uv,x}=r_{u,x}+\sum_{uy=x}r_{v,y}.
\]
Factoring the multiplier proves the assertion, including repeated
prime factors. Consequently, \eqref{eq:kubert} holds in the precise
prime-row form used here for arbitrary $q$, not just squarefree $q$.

In every row of \eqref{eq:ordinary-relations}, the coefficient of $[0]$
is zero. At $x=0$, its occurrences on the two sides cancel; at $x\ne0$,
zero is not a preimage. Thus $[0]$ is a free direct summand of $U(q)$.
Deleting this coordinate leaves a free quotient of rank $\varphi(q)-1$.
Since the quotient is free over $\Z$, reduction modulo any prime
preserves its dimension. In particular, the distribution matrix on the
$q-1$ nonzero coordinates has rank
\begin{equation}\label{eq:ordinary-rank}
 q-\varphi(q)
\end{equation}
over every field $\F_\ell$.

\section{A saturated presentation of the digamma lattice}

We first record the elementary local algebra that will transfer
relations from $\Q$ to $\Zloc$.

\begin{lemma}[Equal-rank saturation]\label{lem:saturation}
Let $B$ be an integer matrix with $m$ columns. Suppose
\[
 \rank_{\F_\ell}(B\bmod\ell)=\rank_\Q B=r.
\]
Let $W_\Q$ be its rational row span and let $W_{\ell}$ be its row
span over $\Zloc$. Then
\begin{equation}\label{eq:saturation}
 W_\Q\cap\Zloc^m=W_{\ell}.
\end{equation}
\end{lemma}

\begin{proof}
Choose $r$ rows and $r$ columns whose minor is nonzero modulo $\ell$.
The chosen rows form a basis of $W_\Q$, and the chosen square minor
is invertible over $\Zloc$. If $v\in W_\Q\cap\Zloc^m$, express $v$
in this rational row basis. Its chosen coordinates determine the
coefficients by multiplication with the inverse minor. All these
coefficients belong to $\Zloc$, proving \eqref{eq:saturation}.
The reverse inclusion is immediate.
\end{proof}

Fix $q$ and $\ell$ satisfying \eqref{eq:residue-hypothesis}. Let
$P=\{p:p\mid q\}$ and $w=|P|$. Introduce the free $\Zloc$-module
\[
 F=\bigoplus_{a=1}^{q-1}\Zloc E_a
   \ \oplus\ \bigoplus_{p\in P}\Zloc Z_p,
 \qquad E_0:=0.
\]
The symbols $Z_p$ are formal logarithm coordinates, not numbers. Define
evaluation by
\begin{equation}\label{eq:evaluation}
 \operatorname{ev}(E_a)=D_a(q),\qquad
 \operatorname{ev}(Z_p)=\log p.
\end{equation}
Equation~\eqref{eq:prime-log-in-span} shows that its image is exactly
$\Lambda_{q,\ell}$.

For $p\mid q$, set $m=q/p$. For every integer $0\le a<m$, define
\begin{equation}\label{eq:augmented-row}
 b_{p,a}=\sum_{j=0}^{p-1}E_{a+jm}-pE_{pa}+pZ_p.
\end{equation}
All displayed indices are between $0$ and $q-1$. Let $B$ be the
integer matrix consisting of these rows, and let $W_\ell$ be its
$\Zloc$-row module.

Every row evaluates to zero. For $a>0$, this is
\eqref{eq:multiplication} with $z=a/q$; for $a=0$, it is
\eqref{eq:prime-log-in-span}. In particular, the zero-index row is
\begin{equation}\label{eq:zero-row}
 b_{p,0}=\sum_{j=1}^{p-1}E_{jq/p}+pZ_p.
\end{equation}
This row is an actual digamma identity, not an evaluation of $\psi(0)$.

\begin{proposition}[Exact presentation and logarithm independence]
\label{prop:presentation}
The map \eqref{eq:evaluation} induces an isomorphism of
$\Zloc$-modules
\begin{equation}\label{eq:module-isomorphism}
 F/W_\ell\ \simeq\ \Lambda_{q,\ell}.
\end{equation}
Moreover, the classes of $Z_p$, $p\mid q$, are linearly independent
in $(F/W_\ell)/\ell(F/W_\ell)$.
\end{proposition}

\begin{proof}
There are $q-1+w$ columns. By \eqref{eq:digamma-dimension}, and because
the logarithms are already in the digamma span, the rational evaluation
map has kernel dimension
\begin{equation}\label{eq:kernel-dimension}
 (q-1+w)-(\varphi(q)+w-1)=q-\varphi(q).
\end{equation}
Its kernel contains every row of $B$, so
$\rank_\Q B\le q-\varphi(q)$.

Reduce $B$ modulo $\ell$ and project onto its $E$ coordinates. Since
$p\equiv1\pmod\ell$, the projected rows are
\[
 \sum_{j=0}^{p-1}E_{a+jm}-E_{pa}.
\]
These are exactly \eqref{eq:ordinary-relations}, with
$x=pa/q$ and $y=(a+jm)/q$, after deleting the zero coordinate.
Their rank is $q-\varphi(q)$ by \eqref{eq:ordinary-rank}. We therefore
have the rank comparison
\begin{align}
 q-\varphi(q)
 &=\rank_{\F_\ell}(\text{$E$-projection of }B)\notag\\
 &\le\rank_{\F_\ell}B
 \le\rank_\Q B
 \le q-\varphi(q).\label{eq:rank-comparison}
\end{align}
All ranks are equal. Equality with \eqref{eq:kernel-dimension} shows
that the rational row span is the entire rational evaluation kernel.
Lemma~\ref{lem:saturation} then shows that every relation with
$\Zloc$ coefficients belongs to $W_\ell$. This proves
\eqref{eq:module-isomorphism}.

Let $\overline W$ be the row span of $B$ over $\F_\ell$. Its
projection to the $E$ coordinates has the same dimension as
$\overline W$, by \eqref{eq:rank-comparison}; hence this projection
is injective. Thus $\overline W$ contains no nonzero vector supported
only on the $Z$ coordinates. A relation among the classes of $Z_p$
in the stated quotient would give precisely such a vector, proving
their independence.
\end{proof}

\begin{remark}
The rank comparison proves both completeness of the rational
distribution relations and their saturation at $\ell$. Neither is
assumed. Integral freeness of the ordinary distribution supplies the
lower bound in characteristic $\ell$; the analytic digamma dimension
supplies the upper bound in characteristic zero.
\end{remark}

\section{Prime-logarithm integrality and power obstructions}

\begin{proof}[Proof of Theorem~\ref{thm:purity-main}]
Each $\log p$ belongs to $\Lambda_{q,\ell}$ by
\eqref{eq:prime-log-in-span}. Proposition~\ref{prop:presentation}
therefore proves the asserted independence in
$\Lambda_{q,\ell}/\ell\Lambda_{q,\ell}$.

Suppose $\sum_{p\mid q}c_p\log p\in\Lambda_{q,\ell}$ with
$c_p\in\Q$. If some coefficient is not in $\Zloc$, choose the
smallest $k\ge1$ such that every $\ell^kc_p$ belongs to $\Zloc$.
At least one of these scaled coefficients is a unit modulo $\ell$.
Multiplying the assumed lattice expression by $\ell^k$ gives
\[
 \sum_{p\mid q}\ell^kc_p\log p
 \ \in\ \ell\Lambda_{q,\ell}.
\]
This contradicts the independence just proved. Hence all
$c_p\in\Zloc$, establishing \eqref{eq:purity-main}. The sum on the
right is direct because the logarithms of distinct primes are
rationally independent, by unique factorization.
\end{proof}

The same lattice identity yields a necessary condition at every power
of the auxiliary prime.

\begin{theorem}[Higher-power obstruction]\label{thm:higher-power}
Assume \eqref{eq:residue-hypothesis} and let $k\ge1$. Suppose $f$ is
$q$-periodic, $f(0)=0$, and
\[
 f(a)\in1+\ell^k\Zloc\quad(1\le a<q),\qquad
 \sum_{a=1}^{q}f(a)=0,\qquad L(1,f)=0.
\]
Then $\ell^k\mid v_p(q)$ for every $p\mid q$. Equivalently, $q$
is a perfect $\ell^k$th power.
\end{theorem}

\begin{proof}
Write $f(a)=1+\ell^kg_a$ with $g_a\in\Zloc$. The zero-mean condition,
$f(0)=0$, and \eqref{eq:digamma-L} give
\[
 0=\sum_{a=1}^{q-1}f(a)D_a(q)
   =-q\log q+\ell^k\sum_{a=1}^{q-1}g_aD_a(q).
\]
Consequently,
\begin{equation}\label{eq:logq-in-lattice}
 \frac q{\ell^k}\log q
 =\sum_{p\mid q}\frac{qv_p(q)}{\ell^k}\log p
 \ \in\ \Lambda_{q,\ell}.
\end{equation}
Theorem~\ref{thm:purity-main} forces every coefficient
$qv_p(q)/\ell^k$ into $\Zloc$. Since $\ell\nmid q$, this is exactly
$\ell^k\mid v_p(q)$.
\end{proof}

The case $k=1$ proves the necessary direction of
Theorem~\ref{thm:classification}. For $k>1$, the theorem is a necessary
condition only; no converse at that precision is asserted here.

\section{Explicit witnesses for perfect powers}

We now prove the converse of Theorem~\ref{thm:classification} and keep
track of integer coefficients and the congruence condition.

\begin{lemma}[Prime-power witness]\label{lem:primepower-witness}
Let $p\equiv1\pmod\ell$ be prime and let $e\ge1$ satisfy $\ell\mid e$.
There is a mean-zero, $p^e$-periodic integer function $f$ with $f(0)=0$,
with $f(a)\equiv1\pmod\ell$ for $1\le a<p^e$, and with $L(1,f)=0$.
\end{lemma}

\begin{proof}
The cyclic group $\F_p^\times$ has a subgroup $K$ of index $\ell$.
Define a $p$-periodic function $h$ by
\begin{equation}\label{eq:h-definition}
 h(a)=
 \begin{cases}
 1-\ell\mathbf1_K(a\bmod p),&p\nmid a,\\
 0,&p\mid a.
 \end{cases}
\end{equation}
Its mean is zero because $|K|=(p-1)/\ell$, and every unit value is
$1$ modulo $\ell$.

Set
\begin{equation}\label{eq:c-definition}
 c_j=1\quad(1\le j<e),\qquad
 c_0=-\sum_{j=1}^{e-1}p^{-j}.
\end{equation}
These coefficients belong to $\Zloc$. Since $p\equiv1\pmod\ell$
and $e\equiv0\pmod\ell$,
\[
 c_0\equiv-(e-1)\equiv1\pmod\ell.
\]
Define a rational-valued function $f_0$ by
\begin{equation}\label{eq:valuation-witness}
 f_0(n)=
 \begin{cases}
 c_jh(n/p^j),&j=v_p(n)<e,\\
 0,&p^e\mid n.
 \end{cases}
\end{equation}
It is $p^e$-periodic. Within the layer $v_p(n)=j<e$, each nonzero
residue modulo $p$ occurs $p^{e-j-1}$ times after division by $p^j$.
Hence every layer has sum zero, so $f_0$ has mean zero. All its nonzero
residue values are $1$ modulo $\ell$ in $\Zloc$.

Partitioning the series into its finitely many valuation layers gives
\[
 L(1,f_0)
 =\sum_{j=0}^{e-1}c_jp^{-j}
   \sum_{\substack{a\ge1\\p\nmid a}}\frac{h(a)}a
 =\left(\sum_{j=0}^{e-1}c_jp^{-j}\right)L(1,h)=0,
\]
where \eqref{eq:c-definition} makes the coefficient zero. Each series
used here converges, since $h$ has mean zero and is periodic.

Finally, multiply $f_0$ by $p^{e-1}$. This clears every denominator in
\eqref{eq:c-definition}. The multiplier is $1$ modulo $\ell$, so
$f=p^{e-1}f_0$ retains the congruence condition, the zero mean, and
the vanishing harmonic sum.
\end{proof}

\begin{lemma}[Composition of periods]\label{lem:composition}
Suppose $f$ and $g$ satisfy the integer-valued conditions of
Theorem~\ref{thm:classification} at periods $u$ and $v$, respectively.
Then such a function exists at period $uv$.
\end{lemma}

\begin{proof}
Define
\[
 F(n)=
 \begin{cases}
 f(n),&u\nmid n,\\
 g(n/u),&u\mid n.
 \end{cases}
\]
This is $uv$-periodic, is zero on multiples of $uv$, and is $1$ modulo
$\ell$ at all other residue classes. Its sum over a period is
\[
 v\sum_{a=1}^{u}f(a)+\sum_{b=1}^{v}g(b)=0.
\]
Because $f$ vanishes on multiples of $u$, splitting finite partial sums
and then passing to the limit gives
\[
 L(1,F)=L(1,f)+\frac1uL(1,g)=0.
\]
No coprimality assumption on $u$ and $v$ is needed.
\end{proof}

\begin{proof}[Completion of Theorem~\ref{thm:classification}]
The implication \textup{(i)}$\Rightarrow$\textup{(iii)} is
Theorem~\ref{thm:higher-power} with $k=1$. If $q$ is a perfect
$\ell$th power, every prime-power factor $p^{v_p(q)}$ has exponent
divisible by $\ell$. Lemma~\ref{lem:primepower-witness} supplies
integer witnesses for these factors, and
Lemma~\ref{lem:composition} combines them. This proves
\textup{(iii)}$\Rightarrow$\textup{(ii)}. The implication
\textup{(ii)}$\Rightarrow$\textup{(i)} is immediate.
\end{proof}

\begin{example}[A witness at period $9$]\label{ex:nine}
For an odd prime $p$, let $\chi$ be the Legendre symbol modulo $p$.
A particularly short integer witness at period $p^2$ is
\[
 f_p(n)=
 \begin{cases}
 -\chi(n),&p\nmid n,\\
 p\chi(a),&n=pa,\ p\nmid a,\\
 0,&p^2\mid n.
 \end{cases}
\]
Both nonzero layers have sum zero, and
$L(1,f_p)=-L(1,\chi)+L(1,\chi)=0$.
For $p=3$, its values on residues $1,\ldots,8$ are
\[
 (-1,\ 1,\ 3,\ -1,\ 1,\ -3,\ -1,\ 1).
\]
This verifies square-period existence with odd integer values while
displaying explicitly why the construction is not a sign-function
construction.
\end{example}

\begin{example}[A witness modulo $3$ at period $343$]\label{ex:343}
Take $\ell=3$, $p=7$, and $e=3$. The subgroup
$K=\{1,6\}\subset\F_7^\times$ has index $3$. Thus $h=-2$ on $K$
and $h=1$ on the other units. The coefficients in
\eqref{eq:c-definition} are
\[
 (c_0,c_1,c_2)=\left(-\frac8{49},1,1\right),
 \qquad c_0+\frac{c_1}{7}+\frac{c_2}{49}=0.
\]
After multiplication by $49$, the integer witness is
\[
 f(n)=
 \begin{cases}
 -8h(n),&7\nmid n,\\
 49h(n/7),&v_7(n)=1,\\
 49h(n/49),&v_7(n)=2,\\
 0,&343\mid n.
 \end{cases}
\]
Every nonzero residue value is $1$ modulo $3$. Its valuation-layer
means vanish, and its harmonic sum is
$(-8+7+1)L(1,h)=0$.
\end{example}

\section{Erd\H{o}s sums and the scope of the conclusion}

\begin{proof}[Proof of Corollaries~\ref{cor:squares} and~\ref{cor:erdos}]
Corollary~\ref{cor:squares} is Theorem~\ref{thm:classification} with
$\ell=2$. If $q$ is odd and nonsquare, a sign function satisfying the
mean-zero condition cannot have $L(1,f)=0$, by this classification.

For completeness, a mean-zero rational periodic function has a finite
Fourier expansion with algebraic coefficients. Abel summation of the
power series gives an expression for $L(1,f)$ as an algebraic linear
combination of logarithms of nonzero algebraic numbers. Baker's
theorem implies that such a linear form is either zero or
transcendental; this is the standard arithmetic alternative used for
periodic Dirichlet series \cite{BBW,MurtyPathak}. Since zero has been
excluded, $L(1,f)$ is transcendental.
\end{proof}

For example, the theorem applies to every mean-zero sign pattern at
period $q=105$. This period is $1$ modulo $4$, and
$2\varphi(105)+1=97<105$, so neither the $3$ modulo $4$ result nor
the elementary totient criterion gives this particular conclusion.
More generally, the present argument treats all nonsquare odd periods
at once. Earlier work includes the $3$ modulo $4$ case and a
density result for periods congruent to $1$ modulo $4$; see the
discussion and references in \cite{ChatterjeeMurty}.

\begin{corollary}[Density of exceptional sign periods]\label{cor:density}
For $X\ge1$, let $\mathcal E(X)$ be the set of integers $2\le q\le X$
for which a $q$-periodic sign function $f$, with $f(0)=0$, has a
convergent and vanishing harmonic sum. Then
\[
 |\mathcal E(X)|\le
 \left\lfloor\frac{\sqrt X-1}{2}\right\rfloor.
\]
Thus the exceptional periods have natural density zero. Among odd
periods, uniform nonvanishing over all convergent sign patterns holds
on a set of natural relative density one.
\end{corollary}

\begin{proof}
An even period cannot have a mean-zero sign pattern with $f(0)=0$,
because the sum of its $q-1$ signs is odd. For odd periods,
Corollary~\ref{cor:erdos} confines every exception to an odd square
$q\ge9$. The displayed bound counts those squares. Dividing by the
number of all periods, or by the number of odd periods, proves the
density assertions.
\end{proof}

The density in Corollary~\ref{cor:density} is a density of moduli; the
assertion at each included modulus is simultaneous for all convergent
sign patterns.

The square obstruction is stronger than a statement about sign
functions: it applies to every odd-valued rational function with
odd denominators. Its converse, however, belongs to that larger
class. Classifying the square periods that admit vanishing
$\{-1,1\}$-valued sums is a separate problem. Likewise,
Theorem~\ref{thm:higher-power} provides only a necessary condition
at precision $\ell^k$ when $k>1$.

The arguments are uniform in $q$. They use no enumeration of periods,
no experimental nonvanishing assumption, and no unproved assertion
about Euler's constant. The essential integral input is the freeness
of the finite-level ordinary distribution; retaining it, rather than
only its rational rank, is what allows the parity and higher-prime
congruence conclusions.

\paragraph{Consistency checks.}
Exact finite calculations checked the augmented and projected ranks
and the square-period membership criterion for all $150$ odd integers
$3\le q\le301$. Further checks covered $71$ admissible pairs
$(\ell,q)$ with $\ell\in\{3,5,7\}$ and $q\le350$. Integer witnesses
at periods $9$, $625$, $343$, and $2197$ passed exact zero-mean,
coefficient-congruence, and vanishing-coefficient checks. These finite
checks are supplementary; the proofs above establish the assertions
for every permitted period.

\begin{thebibliography}{9}

\bibitem{Anderson}
G.~W. Anderson,
\emph{A double complex for computing the sign-cohomology of the universal
ordinary distribution}, preprint (1997),
\href{https://arxiv.org/abs/math/9710225}{arXiv:math/9710225}.
See the finite-level definition in \S3.6, Theorem~4.3.1(I), and \S5.3.

\bibitem{BBW}
A.~Baker, B.~J. Birch, and E.~Wirsing,
\emph{On a problem of Chowla},
Journal of Number Theory \textbf{5} (1973), 224--236.

\bibitem{ChatterjeeMurty}
T.~Chatterjee and M.~R. Murty,
\emph{On a conjecture of Erd\H{o}s and certain Dirichlet series},
Pacific Journal of Mathematics \textbf{275} (2015), 103--113.
\href{https://msp.org/pjm/2015/275-1/pjm-v275-n1-p05-p.pdf}{Publisher's text}.

\bibitem{Kubert}
D.~S. Kubert,
\emph{The universal ordinary distribution},
Bulletin de la Soci\'et\'e Math\'ematique de France
\textbf{107} (1979), 179--202.
\href{https://doi.org/10.24033/bsmf.1891}{doi:10.24033/bsmf.1891}.

\bibitem{MurtyPathak}
M.~R. Murty and S.~Pathak,
\emph{The Okada space and vanishing of $L(1,f)$},
Functiones et Approximatio Commentarii Mathematici
\textbf{66} (2022), 35--57.
\href{https://www.cmi.ac.in/~siddhi/Okada_space_RM_SP_Final.pdf}{Authors' manuscript}.

\bibitem{VMRC055}
Keio,
\emph{Erd\H{o}s's conjecture on $\sum f(n)/n$ is false:
a $2$-adic shadow of Okada's criterion and the least counterexample
$q=105^2$}, research manuscript, 24 October 2026,
Verifiable Mathematical Reasoning Corpus, no.~055.
\href{https://github.com/QEEDDD/Verifiable-Mathematical-Reasoning-Corpus/blob/main/\%5B055\%5D\%20ERDOS\%20CONJECTURE.tex}{Source manuscript}.

\end{thebibliography}
\end{document}
